\begin{frame}{자주 묻는 질문 2·3: 변환이 없으면?}
  \kq{softmax 없이 내적값을 바로 가중치로 쓰면 안 되나요?}
  내적값은 \textbf{범위가 정해져 있지 않고}(음수 가능), 단어마다
  \textbf{스케일이 달라져} ``합이 1'' 보장이 없으니 출력 스케일이
  불안정. softmax는 Week 2.3 로지스틱회귀와 \textbf{동일한 역할} —
  임의의 실수를 ``\textbf{합이 1인, 항상 양수인} 분포''로 바꿔
  안정적인 \textbf{보간(interpolation)}으로.
  \vskip0.4em
  \kq{$V$ 도 같은 임베딩에서 나오는데, 왜 \emph{서로 다른}
  $W_V$ 가 필요한가? $V = X$ 로 못 하나?}
  \begin{itemize}
    \item 가능하지만 \textbf{표현력이 제한}: $V = X$ 라면 출력은
      ``입력 임베딩들의 가중합''이므로 \textbf{한 층의 출력이 항상
      \emph{입력 임베딩 공간} 안에 갇혀} 있다.
    \item $W_V$ 를 두면 ``\textbf{비교에 쓰는 표현}(Q,K)''과
      ``\textbf{전달하는 내용}(V)''을 \emph{서로 독립적으로}
      학습 — 예: ``그것은''의 \emph{Key}는 ``나를 가리킬 명사''에
      높은 점수를 줘야 하고(비교용), \emph{Value}는 ``지쳐 있음''
      같은 \emph{내용}을 담아줘야 — 두 \emph{역할}을 \emph{같은}
      벡터로 동시에 충당하기 어렵다.
  \end{itemize}
  {\footnotesize 참고: ``$Q=K=V=X$''만으로도 \emph{어텐션
  가중치} 자체는 \emph{일단} 계산된다 (앞 실습이 그 설정) —
  가중치 패턴이 \emph{의미적으로 잘} 나오는지는 $W_Q, W_K, W_V$ 가
  \emph{학습}되었을 때의 문제.}
\end{frame}
